% Transcription — fonds Alexandre Grothendieck, shelfmark 46, batch 1 (pages 1-20).
% Established from the facsimile archives/batches/46/batch-01.pdf (cut from a
% copy of 46.pdf fetched through the project relay,
% https://grothendieck-relay.onrender.com/source/46.pdf; PDF page k =
% archivists' page k, no cover sheet), per the skill transcribe-grothendieck.
% The folder is ninety-four pages in five batches; this is the first
% (pages 21-94 are batches 2-5, done separately).
%
% Pass: Opus 5.5 (claude-opus-5-5), 2026-09-26 — first pass, unchecked against the pages by a human.
%
% Pages skipped: none silently. Three leaves are covers inscribed in his hand
% and take no \page marker; their words become headings, with a \note:
% page 1 (tan folder, « Schémas en groupes et fibrés principaux (Divers) »),
% page 3 (« p-structures sur fibrés vectoriels et sur torseurs »), page 9
% (« Frobenius »).
%
% Pages summarised: none. The margins of pages 11-14 and 16 are partly
% compressed into \note{} where the diagonal writing does not read.
%
% Typescript: the folder title announces a « tapuscrit (s.d.) », but no typed
% leaf occurs in pages 1-20; every page here is manuscript in his hand. The
% question of whose typescript it is falls to the later batches.
%
% His pagination: none on these leaves.
%
% What the batch is about. Page 2: a torn leaf carrying a plan of chapters
% I-XIII (plus an unnumbered one) of what is recognisably the EGA
% programme — Langage des schémas, ..., Procédés de construction (descente),
% VII « Schémas en groupes, fibrés principaux », Picard and abelian
% varieties, torsors and cotorsors, intersection theory and RR, étale
% cohomology and fundamental group, resolution of singularities — heavily
% reworked in two inks, with a struck heading « Groupe fondamental ».
% Pages 4-8 (« p-structures »): for an endomorphism f of a sheaf of groups G
% on a site, the stack of G-torsors P with P ≃ P^f is a gerbe iff
% u -> f(u)u^{-1} is an epimorphism, and is then the gerbe of N-torsors,
% N = Ker(f, id); Lang's theorem for f = Frobenius on a smooth group over F_q;
% corollary: such torsors over S/F_q are Galois coverings with group G(F_q);
% application to p-structures phi : E^{(p)} -> E on locally free modules,
% which correspond to étale F_p-local systems (E = Phi ⊗ O_S).
% Page 10: a proposition — p : X' -> X locally of finite presentation is
% unramified (resp. étale) iff the relative Frobenius Fr_{X'/X} is
% unramified / a monomorphism (resp. étale / an isomorphism), with proof.
% Pages 11-16 (« Frobeniuseries »): absolute Frobenius, P_S(X) = X ×_S
% (S, Frob(S)), Frob_S(X) and Ver_S(X), the same for modules and algebras,
% Spec, Proj, V(E), P(E), p(L) = L^{⊗p}; base change; iterated Frobenius;
% perfect base; finite base field; for groups, U_nu(G) = Ker Frob^nu_S(G), and
% a cancelled proposition. Page 17: Lemme 1, a universal homeomorphism of
% finite presentation is affine (first proof cancelled, second kept).
% Page 18: Proposition 2, G is of infinitesimal order nu iff affine with
% augmentation ideal killed by p^nu-th powers; U_nu(G) as an infinitesimal
% neighbourhood of the unit section. Pages 19-20: p-torsion of Pic(X) for X
% proper normal over a perfect field — pPic(X) = Γ(X, O*/O*^p), the
% logarithmic derivative into closed forms, its injectivity, Cartier's
% operator C (image = forms with Cω = ω), and the reduction to a local
% question; breaks off.
%
% Notation fixed once for the batch: the endomorphism of pages 4-6 is his
% gothic f, set \mathfrak{f}; the site \underline{S} underlined as written.
% His Frobenius letters (a double-struck F with r) are set \mathrm{Fr} and
% \mathrm{Frob}; his P_S(X), P_S(E) (a looped capital P) are set P_S; X^{(p/X)}
% as written. Sheaves underlined in his hand (\underline{O}, \underline{\Omega},
% \underline{V}, \underline{P}) keep their underline inside math.
%
% Readings to check first: the heading of page 4 and the adjective used for
% an invertible p-structure on page 7 (both left \ill); the plan of page 2
% (chapters V, VIII, XI, XII); the interlinear insertions on page 4; the
% relation b^{-1}a = f(u)u^{-1} on page 5 (as written; a note flags it); the
% reference « 17.4.10 » on page 10 and « (16.1.8) » at its foot; the prose of
% pages 12 and 14-16; the margins of pages 12, 14 and 16.

\documentclass[11pt,a4paper]{article}
% Le préambule commun à toutes les transcriptions du fonds Grothendieck.
%
% Il tient en une idée : **l'incertitude se compose, elle ne se résout pas.**
% Une transcription qui lisse un mot douteux a détruit la seule chose pour
% laquelle on la faisait — un lecteur doit pouvoir savoir, à chaque endroit,
% ce qui était sur le feuillet et ce qui a été ajouté. Les six macros qui
% suivent sont donc l'appareil critique, et non de la décoration.
%
% Les mêmes commandes sont reconnues par `scripts/render.mjs`, qui produit la
% vue de lecture du site. Une macro ajoutée ici doit l'être aussi là-bas,
% faute de quoi le rendu échouera bruyamment — ce qui est voulu : un rendu
% silencieusement faux est pire qu'un rendu absent.

\ProvidesPackage{grothendieck}[2026/08/08 Transcriptions du fonds Grothendieck]

\usepackage{amsmath,amssymb,amsthm}
\usepackage{mathrsfs}
\usepackage{stmaryrd}
% Les diagrammes commutatifs sont partout dans ce fonds : c'est la langue
% même de la géométrie algébrique de Grothendieck, et une transcription qui
% les rendrait en prose ne transcrirait rien.
\usepackage{tikz-cd}
% Français par défaut — c'est la langue des feuillets — mais l'anglais est
% chargé aussi : la traduction et le résumé sont des documents anglais, et la
% typographie française (espaces avant les deux-points) y serait une faute.
% Ces documents basculent par \selectlanguage{english} en tête de corps.
\usepackage[english,french]{babel}
\usepackage{fontspec}
\usepackage{geometry}
\geometry{a4paper,margin=25mm,marginparwidth=28mm}
\usepackage{marginnote}
\usepackage{xcolor}
% ulem, not soul. `soul`'s \st and \ul are built on a character-by-character
% scan that XeLaTeX's font machinery breaks: the document fails with an
% undefined control sequence rather than degrading. ulem does the same two jobs
% and is Unicode-engine safe. `normalem` stops it hijacking \emph, which the
% transcriptions use for Grothendieck's own emphasis.
\usepackage[normalem]{ulem}
\usepackage{hyperref}
\hypersetup{colorlinks=true,linkcolor=black,urlcolor=black,pdfborder={0 0 0}}

% --- Métadonnées du lot -----------------------------------------------------
% Reprises telles quelles par le script de rendu, qui les affiche en tête de
% la vue de lecture. Elles fixent ce que le fichier prétend couvrir : sans
% elles, une transcription flotte sans point d'attache dans un fonds de seize
% mille feuillets.

\newcommand{\folder}[1]{\gdef\grothFolder{#1}}
\newcommand{\batch}[1]{\gdef\grothBatch{#1}}
\newcommand{\leaves}[2]{\gdef\grothFirst{#1}\gdef\grothLast{#2}}
\newcommand{\foldertitle}[1]{\gdef\grothFoldertitle{#1}}
\gdef\grothFolder{?}\gdef\grothBatch{?}\gdef\grothFirst{?}\gdef\grothLast{?}\gdef\grothFoldertitle{}

% --- Le résumé --------------------------------------------------------------
%
% Il ouvre la lecture modernisée, sous le titre « Résumé », et il est composé
% en petit. Ce n'est pas un ornement : c'est le seul passage du document qui
% ne soit pas des mathématiques, et un lecteur doit pouvoir le reconnaître
% avant de l'avoir lu — puis le sauter s'il connaît déjà le sujet, ou s'y
% arrêter s'il ne le connaît pas. Le filet qui le suit marque la reprise du
% corps.

\newenvironment{resume}
  {\small\setlength{\parskip}{0.45em}}
  {\par\medskip\hrule\medskip}

% --- Mots-clés ---------------------------------------------------------------
%
% En anglais, en clôture du résumé de la lecture modernisée : le vocabulaire
% moderne sous lequel chercher ce que les feuillets construisent. Le manifeste
% du site (`npm run manifest`) les extrait pour étiqueter la cote entière —
% cette ligne est la source unique des tags, il n'y a pas de fichier à côté.
\newcommand{\keywords}[1]{\par\smallskip\noindent
  {\footnotesize\textsc{Keywords} --- \textit{#1}}\par}

% --- L'appareil critique ----------------------------------------------------

% Le numéro de feuillet, en marge. C'est l'ancre qui permet de régler un
% désaccord sur une formule en regardant *un* feuillet plutôt qu'une chemise
% de six cents. Le site s'en sert aussi pour tourner les pages du fac-similé
% quand on fait défiler la transcription.
\newcommand{\leaf}[1]{%
  \marginnote{\footnotesize\color{gray!70}\textsf{#1}}%
  \label{leaf:#1}\ignorespaces}

% Illisible : on ne devine pas. Un mot inventé qui se lit comme les autres est
% le pire résultat possible.
\newcommand{\ill}{\textcolor{red!60!black}{[\,\ldots\,]}}

% Lecture douteuse : le mot est donné, et signalé comme incertain.
\newcommand{\uncertain}[1]{\uline{#1}}

% Ajout de l'éditeur — une lettre manquante, un mot sous-entendu.
\newcommand{\add}[1]{\textcolor{blue!50!black}{[#1]}}

% Biffé par Grothendieck lui-même. Ce qu'il a rayé fait partie du document :
% une rature dit souvent où la pensée a changé de direction.
\newcommand{\struck}[1]{\sout{#1}}

% Note du transcripteur, sur l'état matériel du feuillet.
\newcommand{\note}[1]{\par\smallskip{\small\color{gray!60!black}\textit{[#1]}}\par\smallskip}

% Note marginale de Grothendieck, distincte du corps du texte.
\newcommand{\marginal}[1]{\par\smallskip\noindent\hspace{1em}%
  {\small\color{gray!60!black}#1}\par\smallskip}

% --- Titre ------------------------------------------------------------------

\AtBeginDocument{%
  \begin{center}
    {\large\bfseries Fonds Alexandre Grothendieck --- cote n\textsuperscript{o}~\grothFolder}\\[2pt]
    {\itshape\grothFoldertitle}\\[4pt]
    {\small feuillets \grothFirst--\grothLast\ (lot \grothBatch)}\\[2pt]
    {\footnotesize\color{gray!60!black}Transcription établie d'après le fac-similé numérisé
     par l'Université de Montpellier.\\
     Les lectures incertaines sont soulignées, les illisibles notées [\,\ldots\,],
     les ajouts entre crochets.}
  \end{center}
  \vspace{1em}}

\endinput


\folder{46}
\batch{1}
\pages{1}{20}
\dating{[vers 1966-vers 1972]}
\watermark{Édition de démonstration}
\foldertitle{Schémas en groupes et fibrés principau[x] (Divers) : notes manuscrites (s.d.), tapuscrit (s.d.)}

\begin{document}

\section*{Schémas en groupes et fibrés principaux (Divers)}

\note{inscrit à l'encre en haut à droite de la chemise (page 1), qui ne porte
rien d'autre ; le feuillet ne reçoit pas de numéro de page}

\page{2}

\marginal{\struck{Groupe fondamental}}\note{en haut à droite, sous la
déchirure du feuillet, barré de traits obliques ; lecture de « fondamental »
incertaine}

\begin{itemize}
  \item[I] Langage des schémas
  \item[II] Étude élémentaire de quelques classes de morphismes
  \item[III] Cohomologie des faisceaux cohérents.
  \item[IV] Étude locale des morphismes
  \item[V] Étude globale élémentaire des \struck{quelques} morphismes.
  \item[VI] Procédés de construction. \uncertain{d}
  \marginal{[N.B. y inclure propriétés de descente pour morph. étales]}
  \item[VII] Schémas en groupes, fibrés principaux [définition du gp
  fondamental]
  \item[VIII] \struck{Étude cohomologique des faisceaux \ill{} en
  \uncertain{jours} (\ill{} \ill{} \ill{} \ill{} \ill{} ; \ill{} \ill{})}
  \add{Et Picard et V.A.}\note{« Et Picard et V.A. » est écrit à l'encre
  bleue au-dessus de la ligne barrée, dont la lecture est très douteuse}
  \item[\struck{IX}] \struck{Groupe fondamental}/K \uncertain{pour} \ill{}
  la top. étale\ldots
  \item[\struck{X}] \struck{\ill{}}
  \item[IX] Torseurs et cotorseurs, dualité
  \item[X] Théorie d'intersection. RR.
  \item[XI] \struck{Topologie : \ill{} \ill{}, \ill{}} (Cohomologie
  \uncertain{étale} et groupe fondamental)
  \item[XII] \uncertain{Jacobiennes et variétés} \ill{}, \ill{} \uncertain{etc}.
  \item[XIII] Résolution des singularités ?
  \item[\ldots] Théorie des cycles algébriques ??
\end{itemize}

\note{plan de chapitres, écrit à l'encre noire et repris à l'encre bleue :
les numéros ont été récrits (VIII sur VIII barré, IX et X barrés puis
réattribués, XI récrit, XII et XIII récrits) ; une accolade réunit IX et X
avec en regard « intersectivité » \uncertain{}, et dans la marge gauche
« intersectiviste » \uncertain{} avec des flèches déplaçant les chapitres IX à
XI ; un crochet et un « ? » entourent XII et XIII. L'ordre et la numérotation
sont donnés dans leur dernier état lisible}

\section*{p-structures sur fibrés vectoriels et sur torseurs}

\note{inscrit en haut à droite d'un feuillet par ailleurs blanc (page 3), qui
ne reçoit pas de numéro de page}

\page{4}

\emph{p-structures \ill{} \ill{}.}

1) Soit $\underline{S}$ un site, $G$ un faisceau de groupes sur
$\underline{S}$,
\[
  \mathfrak{f} : G \longrightarrow G
\]
un endomorphisme. On s'intéresse à la structure suivante (où, pour un
torseur $P$ sous $G$, $P^{\mathfrak{f}}$ désigne le torseur déduit par
extension du groupe par $\mathfrak{f}$) : donnée d'un torseur $P$ sous $G$,
et d'un \emph{isom.} $P \simeq P^{\mathfrak{f}}$, ou ce qui revient au même,
d'un hom.\ de faisceaux
\[
  \varphi : P \longrightarrow P
\]
satisfaisant
\[
  \varphi(xg) = \varphi(x)\,\mathfrak{f}(g) .
\]
Ces objets (ou encore \uncertain{d'objets} \ill{} de $S$)
\add{\uncertain{sur un groupoïde, avec une action} \ill{} $(G, \mathfrak{f})$}
forment un champ\note{l'insertion interlinéaire est reliée au texte par un
trait ; lecture douteuse}. On se donne des conditions moyennant lesquelles ce
champ est une gerbe. Notons que \emph{loc.}\ un tel $P$ est trivial comme
simple torseur, donc peut s'identifier à $G$ pris comme torseur trivial. La
structure $\varphi$ est alors déterminée par la connaissance de
$\varphi(1) = a$, par la formule
\[
  \varphi_a(g) = a\,\mathfrak{f}(g)
\]

\page{5}

et $a$ peut être pris arbitrairement. On voit déjà que le champ envisagé est
\emph{connexe}, et il revient à \struck{\ill{}} \add{voir} la condition
moyennant laquelle deux objets sont loc.\ isom., i.e.\ moyennant laquelle,
pour deux éléments $a, b \in G(S)$ ($S \in \underline{S}$), $\varphi_a$ et
$\varphi_b$ définissent des objets loc.\ isom. Pour ceci, notons que les hom.\
de $(G, \varphi_a)$ dans $(G, \varphi_b)$ s'identifient aux $u \in G(S)$ tels
que

\begin{tikzcd}
G \arrow[r, "\varphi_a"] \arrow[d, "\rho_u"'] & G \arrow[d, "\rho_u"] \\
G \arrow[r, "\varphi_b"'] & G
\end{tikzcd}

commute, i.e.\ (où $\rho_u(g) = ug$) tels que
\[
  u\,a\,\mathfrak{f}(g) = b\,\mathfrak{f}(ug) \qquad \forall g \in G(S'),\
  S' \text{ sur } S,
\]
ce qui s'écrit aussi simplement
\[
  b^{-1}a = \mathfrak{f}(u)\,u^{-1} .
\]
\note{la formule est telle qu'écrite ; de $ua = b\,\mathfrak{f}(u)$ on tire
$b^{-1}ua = \mathfrak{f}(u)$, qui ne se réduit à la forme écrite que si $u$
commute à $a$. La lettre notée ici $\rho_u$ est une lettre peu lisible sur les
flèches verticales}

Donc on trouve que $\Sigma$ \emph{est une gerbe ssi l'hom.\ de faisceaux
d'ensembles}
\[
  u \longmapsto \mathfrak{f}(u)\,u^{-1} : G \longrightarrow G
\]
\emph{est un épimorphisme}. Alors l'objet de $\Sigma$ est loc.\ isom.\ à
$(G, \varphi_e = \mathfrak{f})$ lui-même, dont le faisceau des autom.\ est
formé \struck{des} du faisceau des $u$ de $G$ tels que
$\mathfrak{f}(u)u^{-1} = 1\cdot 1^{-1} = 1$,

\page{6}

donc le faisceau $N = \operatorname{Ker}(\mathfrak{f}, \mathrm{id})$ des
coïncidences de $\mathfrak{f}$ et de $\mathrm{id}_G$. Donc $\Sigma$
\emph{équivaut à la gerbe des $N$-torseurs}.

(2)\note{le numéro est cerclé} Soit $G$ un schéma en groupes lisse
\struck{connexe} sur le corps fini $\mathbb{F}_q$, $q = p^n$. Soit
$\mathfrak{f} = \mathrm{Frob}^n_{G/k} : G \to G$.

\emph{Théorème} (Lang). \textit{Le morphisme $u \mapsto \mathfrak{f}(u)u^{-1}$
de $G$ dans $G$ est \emph{étale}, et de plus \emph{surjectif} (donc couvrant
pour la top.\ étale) si $G$ est connexe.}

Le fait que le morphisme soit \emph{étale} provient du critère
infinitésimal d'étalité, grâce au fait que \struck{\ill{}}
$\mathrm{Lie}(\mathfrak{f}) = 0$ et grâce au fait que $G$ est lisse (i.e.\
\ill{} \ill{} \ill{}). \struck{[De cette façon]}
[De cette façon, $N = \operatorname{Ker}(\mathfrak{f}, \mathrm{id}_G)$ est un
gr.\ étale \add{fini} sur $k$, et
\[
  G/N \xrightarrow[\ \varphi\ ]{\text{isom}} G
\]
induit par $\mathfrak{f}(u)u^{-1}$.]
\struck{Le morphisme $\varphi$ est \ill{}} Kif-kif pour
\[
  \boxed{\varphi(gx) = \mathfrak{f}(g)\,\varphi(x)\,g^{-1}}
\]
morphismes $u \mapsto \mathfrak{f}(u)u^{-1}$. On termine, si $G$ est connexe,
comme dans Serre \uncertain{Gr.\ alg.\ et corps de classes}, p.~119.

\emph{NB} Axiomatisation : $G$ lisse sur $S$, $\mathfrak{f}$ endom.\ tel que
$\mathrm{Lie}(\mathfrak{f}) = 0$. Alors $u \mapsto \mathfrak{f}(u)u^{-1}$ est
étale, et si $G$ est à fibres connexes, il est surjectif.
\note{l'axiomatisation est ajoutée en bas de page, d'une écriture plus
serrée ; une insertion interlinéaire au-dessus de « Serre », « \ill{} un
\uncertain{gpe} », n'est pas lue}

\page{7}

\emph{Corollaire} Soit $S$ sur $\mathbb{F}_q$. La catégorie des torseurs $P$
sur $S$ sous $G_S$, munis de $P \simeq P^{\mathfrak{f}}$, i.e.\ de
$\varphi : P \to P$ satisfaisant
$\varphi(xg) = \varphi(x)\,\mathfrak{f}(g)$, équivaut à celle des
\struck{revêtements \add{étales} galoisiens de $S$, de groupe}
$G(\mathbb{F}_q)$.
\note{la ligne « revêtements \ldots\ de groupe » est barrée d'un trait
plein puis soulignée d'un pointillé, qui paraît la rétablir ; « étales » est
écrit au-dessus}

En effet, $N$ est le groupe discret \add{constant} (\ill{} \ill{}) défini par
$G(\mathbb{F}_q)$, et on applique le n°~1.

(3)\note{numéro cerclé} Soit $S$ préschéma de car.\ $p$, et $E$ un Module
(qu.\ coh.\ !) sur $S$. Une $p$-structure sur $E$ est un hom.
\[
  \varphi : E^{(p)} \longrightarrow E
\]
[où $E^{(p)} = E \otimes_{\mathcal{O}_S} (\mathcal{O}_S, \mathfrak{f})$]. La
$p$-structure est dite \emph{\ill{}} si $\varphi$ est un isom. On veut
classifier les $p$-structures \ill{} sur $S$, avec $E$ loc.\ libre de rang
$n$. Or ce sont les structures de $G$-torseur avec
$\mathfrak{f}$-\ill{}, où $G = \mathrm{GL}(n)_S$, et
$\mathfrak{f} : G \to G$ le Frobenius. Comme $G$ provient de $G_0$ sur
$\mathbb{F}_p$, on peut appliquer 2° ; on trouve que

\page{8}

la catégorie en question est celle des torseurs sous
$\mathrm{GL}(n, \mathbb{F}_p)$, ou ce qui revient au même, celle des
faisceaux (ét) en \struck{groupes} vectoriels sur $\mathbb{F}_p$, qui sont
loc.\ libres de rang $n$. Si $\Phi$ est un tel faisceau, le Module $E$ qui lui
correspond est simplement
\[
  E = \Phi \otimes_{\mathbb{F}_p} \mathcal{O}_S
\]
[ici on note que les faisceaux qu.\ coh.\ (Zariski) et les faisceaux qu.\
coh.\ étales, c'est pareil\ldots].

On peut se demander dans quelle mesure les hypothèses de loc.\ liberté et de
type fini sont nécessaires dans cet énoncé.

\section*{Frobenius}

\note{inscrit en haut à droite d'un feuillet par ailleurs blanc (page 9), qui
ne reçoit pas de numéro de page}

\page{10}

\begin{tikzcd}[column sep=small, row sep=small]
 & X'^{(p/X)} \arrow[dl] \arrow[ddr] & \\
X' \arrow[d, "p"'] & & X' \arrow[ll, "\mathrm{Fr}_{X'}"'] \arrow[ul, "\mathrm{Fr}_{X'/X}"'] \arrow[d, "p"] \\
X & & X \arrow[ll, "\mathrm{Fr}_X"]
\end{tikzcd}

\note{diagramme en haut à gauche, les flèches de Frobenius dirigées vers la
gauche ; sous le $X$ de gauche : « de car.~$p$ ». Le trait qui descend de
$X'^{(p/X)}$ vers le $X$ de droite est fin et sans tête nette}

\emph{Proposition} \emph{Supposons} $p$ loc.\ de prés.\ finie.

1°) Conditions équivalentes
\begin{itemize}
  \item[(i)] $p$ non ramifié
  \item[(ii)] $\mathrm{Fr}_{X'/X}$ non ramifié
  \item[(ii bis)] $\mathrm{Fr}_{X'/X}$ un monomorphisme
\end{itemize}

2°) Conditions équivalentes
\begin{itemize}
  \item[(i)] $p$ étale
  \item[(ii)] $\mathrm{Fr}_{X'/X}$ étale
  \item[(ii bis)] $\mathrm{Fr}_{X'/X}$ un isomorphisme
\end{itemize}

\emph{Dém} 1°) (ii) $\Longleftrightarrow$ (ii bis) car monomorphisme $=$
radiciel non ramifié (\uncertain{17.4.10}).

(i) $\Longleftrightarrow$ (ii) car
$\underline{\Omega}^1_{X'/X} \simeq \underline{\Omega}^1_{X'/X'^{(p/X)}}$.

2°) (ii) $\Longleftrightarrow$ (ii bis) car isom $=$ radiciel surjectif
étale.

(i) $\Longrightarrow$ (ii) car si $X'$ étale sur $X$, $X'^{(p/X)}$ itou, et
tout $X$-morphisme $X' \to X'^{(p/X)}$ itou.

(ii bis) $\Longrightarrow$ (i). Grâce à 1°) on sait déjà que $X'/X$ est
\emph{non ramifié}, donc localement \add{(ét)} au-dessus de $X$
\struck{\ill{}} c'est une immersion. Donc on est ramené à ceci : si
$X' \to X$ est une immersion de prés.\ finie telle que $X' \to X'^{(p/X)}$
est un isom., alors c'est une imm.\ ouverte. On peut supposer que c'est une
imm.\ fermée définie par Idéal $\mathcal{J}$ de type fini, alors $X'^{(p/X)}$
est défini par l'Idéal $\mathcal{J}^{p}$, l'hyp.\ s'écrit
$\mathcal{J} = \mathcal{J}^{p}$, ce qui implique que \struck{\ill{}}
$X' \to X$ est une immersion ouverte (\uncertain{16.1.8}), O.K.

\page{11}

\emph{Frobeniuseries.}

\marginal{$\mathrm{Fr}_X$}\note{cerclé dans la marge gauche}

On désigne par $(\mathrm{Sch})_p$
$\bigl[= (\mathrm{Sch})_{/\operatorname{Spec}(\mathbb{Z}/p\mathbb{Z})}\bigr]$
la sous-catégorie \struck{pleine} de $(\mathrm{Sch})$, formée des préschémas
de car.\ $p$, et tous les préschémas envisagés sont de car.\ $p$.

On définit un endom.\ fonctoriel de $X$
\[
  \mathrm{Frob}(X) : X \longrightarrow X
\]
[\emph{Frobenius absolu}] par la condition d'être l'identité sur les espaces
topologiques, et $f \mapsto f^p$ sur les fonctions.

Soit $S$ fixé. Pour tout $X$ sur $S$, on a donc un carré commutatif

\begin{tikzcd}
X \arrow[r, "\mathrm{Frob}(X)"] \arrow[d, "f"'] & X \arrow[d, "f"] \\
S \arrow[r, "\mathrm{Frob}(S)"'] & S
\end{tikzcd}

\qquad $X \xrightarrow{\ \mathrm{Fr}_{X/S}\ } X^{(p/S)}$

(où $f$ est le morphisme structural). \struck{Ainsi}, Posons
\[
  P_S(X) = X \times_S (S, \mathrm{Frob}(S)) \qquad
  P_S : (\mathrm{Sch})_{/S} \longrightarrow (\mathrm{Sch})_{/S}
\]
$P_S(X)$ est un foncteur, qui \uncertain{étant} un foncteur de changement de
base, \uncertain{a toutes} les propriétés [commute aux \emph{lim}, aux sommes,
transforme type fini

\marginal{$X^{(p/S)}$ ou $X^{(p/f)}$ ou $X^{(p)}$, $\;= (X/S)^{(p)}$}
\note{dans la marge gauche, entouré d'un ovale, avec au-dessus, de biais,
« \uncertain{Changer} \ill{} \ill{} flèches \ill{} direction », et au-dessous,
de biais et en partie barré, « dire que $X \to P_S(X)$ est un
\uncertain{homéomorphisme universel} »}

\page{12}

resp.\ affine resp.\ \ldots\ en type fini resp.\ affine resp.\ldots].

\marginal{$\mathrm{Fr}_{X/S}$}\note{cerclé dans la marge gauche, avec le
diagramme suivant}

\begin{tikzcd}[column sep=large]
X \arrow[r, "\mathrm{Frob}_S(X)"] \arrow[rr, bend left=30, "\mathrm{Frob}(X)"] & P_S(X) \arrow[r, "\mathrm{pr}_1"] \arrow[d] & X \arrow[d] \\
 & S \arrow[r, "\mathrm{Frob}(S)"'] & S
\end{tikzcd}

Le carré précédent nous donne un $S$-morphisme fonctoriel en $X$
\[
  \mathrm{Frob}_S(X) : X \longrightarrow P_S(X)
\]
[et $\mathrm{Ver}_S(X) : P_S(X) \to X$ compatible avec $\mathrm{Frob}(S)$].

Définitions analogues pour un Module $E$ sur $X$, on pose
\[
  P_S(E) = E \otimes_{\mathcal{O}_X} \mathcal{O}_{P_S(X)}
  = E \otimes_{\mathcal{O}_S} (\mathcal{O}_S, \mathrm{Frob}(S)) ,
\]
\struck{$= E$}. En particulier pour $X = S$, on trouve
\[
  P_S(E) = E \otimes_{\mathcal{O}_S} (\mathcal{O}_S, \mathrm{Frob}(S)) .
\]
C'est un foncteur \ill{} en le $\mathcal{O}_X$-Module $E$.
\struck{$\mathrm{Frob}_S(E) : P_S(E) \to E$} \struck{\ill{}} Il
\uncertain{commute} aussi \ill{} \ill{} $S$ \ill{} \ill{} \ill{}, et \ill{}
\ill{} \ill{}. Il \ill{} \ill{} des propriétés \uncertain{générales} des
\ill{} \ill{} \ill{} \ill{} \ill{} $P_S(\mathcal{A})$ \ill{}, \ill{} \ill{}
\uncertain{qu.\ coh.\ et type fini}, \uncertain{qu.\ coh.\ et présentation
finie}, \ill{} \ill{}) il en est de même de $P_S(E)$. \ill{} \ill{}
\emph{$E \mapsto P_S(E)$ \ill{} \ill{} \ill{}}, \ill{} \ill{} \ill{}
$\otimes$, mais \ill{} \ill{} \ill{} \ill{} \ill{} : \ill{} \ill{} \ill{}
\ill{} \ill{} \ill{}.
\note{la prose de la seconde moitié de la page est très rapide ; seuls les
éléments donnés se lisent. Un passage (« C'est un foncteur \ldots\ Module $E$
\ldots\ ») est encadré d'une accolade à gauche}

\note{marge gauche, écrit de biais : « Rôle universel de
$F^{*}(p) = P_S(E)$ \ill{} » et « \ill{} \uncertain{canoniquement} \ill{}
\ill{} $\mathrm{Frob}_S(X)^{*}(P_S(E)) \simeq p(E)$ »}

\page{13}

Si $E = \mathcal{A}$ est une Algèbre qu.\ coh.\ \add{(resp.\ et gr.\
\ill{})}, il en est de même de $P_S(\mathcal{A})$, et on a aussi :
\[
  \begin{cases}
    P_S(\operatorname{Spec}(\mathcal{A})) = \operatorname{Spec}(P_S(\mathcal{A})) \\
    P_S(\operatorname{Proj}(\mathcal{A})) = \operatorname{Proj}(P_S(\mathcal{A}))
  \end{cases}
\]
et le morphisme $\mathrm{Frob}_S(X)$ pour $X = \operatorname{Spec}(\mathcal{A})$,
resp.\ $X = \operatorname{Proj}(\mathcal{A})$, est induit par
$\mathrm{Frob}_S(\mathcal{A}) : P_S(\mathcal{A}) \to \mathcal{A}$, qui
\ill{} : \ill{} \uncertain{puissance $p$-ième} (\ill{}).

\note{dans la marge gauche, un diagramme :}

\begin{tikzcd}[column sep=small, row sep=small, nodes={font=\scriptsize}]
 & P_S(X') \arrow[r] \arrow[d] & X' \arrow[d] \\
X \arrow[r] \arrow[dr, "f"'] & P_S(X) \arrow[r] \arrow[d] & X \arrow[d, "f"] \\
 & S \arrow[r, "\mathrm{Frob}(S)"'] & S
\end{tikzcd}

Corollaire :
\struck{Exemples : $E = \underline{\mathcal{O}}$}
\[
  \begin{cases}
    P_S(\underline{V}(E)) = \underline{V}(P_S(E)) \\
    P_S(\underline{P}(E)) = \underline{P}(P_S(E))
  \end{cases}
\]
Les morphismes de Frobenius pour les préschémas sont ceux associés à
\struck{$P_S(E)$} \struck{$\mathrm{Frob}_S(E)$~:}
\struck{$P_S(E) \to E$}.
\struck{\ill{} \ill{} \ill{}} \add{\uncertain{l'hom}}
\[
  E^{(p)} \longrightarrow \mathrm{Symm}^{p}(E) \subset \mathrm{Symm}(E)
\]
\uncertain{induisant} \ill{} l'élévation à la puissance $p$-ième.

[N.B. La structure vectorielle de $\underline{V}(E)$ n'est pas
\uncertain{conservée}, mais bien entendu la structure de \ill{} \ill{}
conservée. Notons que dans l'hom.\ de Frobenius,
$\underline{P}(E) \to P_S(\underline{P}(E)) = \underline{P}(P_S(E))$, comme

\begin{tikzcd}[column sep=small, row sep=small, nodes={font=\scriptsize}]
X' \arrow[r, "\mathrm{Frob}_X(X')"] \arrow[rr, bend left=30, "\mathrm{Frob}_S(X')"] \arrow[dr, "g"'] & P_X(X') \arrow[r, "r"] \arrow[d] & P_S(X') \arrow[r] \arrow[d] & X' \arrow[d, "g"] \\
 & X \arrow[r, "\mathrm{Frob}_S(X)"] \arrow[rr, dashed, bend right=30, "\mathrm{Frob}(X)"'] \arrow[dr, "f"'] & P_S(X) \arrow[r] \arrow[d] & X \arrow[d, "f"] \\
 & & S \arrow[r, "\mathrm{Frob}(S)"'] & S
\end{tikzcd}

\note{diagramme en bas à gauche de la page, entouré d'une accolade, avec la
légende « trois carrés cartésiens » ; la flèche $\mathrm{Frob}(X)$ est tracée
en tirets}

\page{14}

un hom.\ de l'Algèbre graduée
\[
  \mathrm{Symm}\,E \longrightarrow (\mathrm{Symm}\,E)_{(p)}
  \qquad (\text{degrés multiples de } p)
\]
le faisceau inversible \ill{} \ill{} de
$\underline{\mathcal{O}}_{\underline{P}(P_S(E))}(1)$ \ill{}
$\underline{\mathcal{O}}_{\underline{P}(E)}(p)$. \ill{} \ill{} \ill{}
\uncertain{particulièrement},

Si $L$ est un Module \emph{inversible}, on a un isom.\ \emph{canonique}
\[
  p(L) \simeq L^{\otimes p} .
\]

\note{dans la marge gauche, écrit verticalement : « Exemple :
$E = \mathcal{O}_S^n$, d'où $P_S(E) \simeq \mathcal{O}_S^n$, donc
$P_S(S[t_1, \ldots, t_n]) \simeq S[t_1, \ldots, t_n]$, l'hom.\ \ill{} étant
donné par $t_i \mapsto t_i^p$. Si \ill{} \ill{} \ill{} \ill{} de $E$ défini
par des équations, \ill{} \ill{} \emph{affine} \ill{} \ill{} \ill{} \ill{}.
Analogue pour des \ill{} \ill{} \ill{} \ill{} \ill{} » ; au début, un
$P_S(\underline{V}(E))$ est barré}

\emph{Changement de base dans Frobenius}

\struck{[cf.\ diagramme]}

1°) (Cf.\ diagramme)
\[
  P_X(X') \simeq P_S(X') \times_{P_S(X)} (X, \mathrm{Frob}_S(X)) ,
\]
et \ill{} \ill{} \ill{} \uncertain{identification}, $\mathrm{Frob}_X(X')$ est
défini \ill{} \ill{} \ill{} \ill{} \uncertain{à partir de}
$\mathrm{Frob}_S(X')$ et $\mathrm{Frob}_S(X)g$.

2°) Si on a un changement de base $S' \to S$, on a
\[
  P_S(X) \times_S S' \simeq P_{S'}(X \times_S S')
\]
Cet isom.\ est \uncertain{fonctoriel} en $X$, \struck{\ill{} \ill{}}
\ill{} \ill{} \ill{} \ill{} \ill{} $\mathrm{Frob}_S(X)$ en
$\mathrm{Frob}_{S'}(X \times_S S')$.

\page{15}

Frobenius itérés, notations $P^{\nu}_S(X)$, $\mathrm{Frob}^{\nu}_S(X)$. Les
développements précédents s'étendent, en remplaçant partout les puissances
$p$-ièmes par les puissances $p^{\nu}$-ièmes.

\emph{Cas d'un préschéma de base parfait}. On dit que $S$ est \emph{parfait}
si \struck{\ill{}}
$\mathcal{O}_S \xrightarrow{\mathrm{Frob}(S)^{*}} \mathcal{O}_S$ est un
\emph{isomorphisme}. Le cas le plus important est celui où $S$ est le
spectre d'un corps, i.e.\ le corps considéré est alors \uncertain{parfait}.

\emph{Dans ce cas, le préschéma $P^{\nu}_S(X)$ est isomorphe au préschéma $X$
par $\mathrm{pr}_1$, mais on se rappellera que ce n'est pas un $S$-isom.
En particulier, $P^{\nu}_S(X)$ n'est pas $S$-isom.\ à $X$.}

\emph{Cas d'un corps de base fini}, à $q = p^{\nu}$ \struck{\ill{}} éléments.
Alors $\mathrm{Frob}^{k\nu}_S(S) : S \to S$ est l'identité. Donc pour de
\emph{tels} $\nu$, \struck{\ill{} \ill{}} $P^{\nu}(X) \to X$ est un
\emph{$S$-isom.} Donc on peut considérer $\mathrm{Frob}^{k\nu}_S(X)$ comme un
endomorphisme

\note{« $\mathrm{Frob}^{k\nu}_S(S)$ » : l'exposant est surchargé, et
l'insertion au-dessus de « $q = p^{\nu}$ » (« \uncertain{éléments},
\ill{} \ill{} ») est en partie raturée}

\page{16}

de $X$, qui n'est d'ailleurs autre que
\[
  \bigl(\mathrm{Frob}^{\nu}_S(X)\bigr)^{k} ,
\]
comme on le constate aussitôt.

D'ailleurs, dans le cas où $\nu = 1$, i.e.\ sur la base absolue
$\mathbb{Z}/p\mathbb{Z}$, on trouve pour
$\mathrm{Frob}_{\mathbb{Z}/p\mathbb{Z}}(X)$ le Frobenius absolu du début.

\emph{Corollaire} Si $S$ est sur $\mathbb{F}_q$, et si $X$ provient d'un
préschéma \struck{$X_0$} \add{$X_0$} sur $\mathbb{F}_q$ par extension de la
base\ldots\note{la phrase s'arrête ; « $X_0$ sur $\mathbb{F}_q$ par extension
de la base » est souligné d'un trait ondulé}

\emph{Cas des groupes}. Si $G$ est muni d'une structure algébrique finie sur
$S$, il en est de même de $P_S(G)$, et \struck{$\mathrm{Frob}_S(X)$}
\[
  \mathrm{Frob}^{\nu}_S(G) : G \longrightarrow P^{\nu}_S(G)
\]
est un morphisme pour les structures en question. Considérons, en
particulier,
\[
  U_{\nu}(G) = \operatorname{Ker}\bigl(\mathrm{Frob}^{\nu}_S(G)\bigr) .
\]

\struck{\emph{Proposition} Si $G$ est de présentation finie sur $S$, alors
$U_{\nu}(G)$ est de présentation finie et fini (surjectif) radiciel sur $S$.}
\note{la proposition est barrée de traits obliques ; « surjectif » est une
insertion interlinéaire}

\note{dans la marge gauche, écrit de biais et relié par un trait à
« Proposition » : « \uncertain{Quid} \ill{} \ill{} \ill{} \ill{} \ill{}
\ill{} $U_{\nu}$ \ill{} \ill{} \ill{} $\mathrm{Frob}^{\nu}$ \ill{}
\uncertain{fini} »}

\page{17}

\emph{Lemme 1} \textit{Un morphisme \struck{radiciel surjectif} \add{de
présentation finie} $f : T \to S$ qui est un homéomorphisme
\struck{universel} est affine.}

\struck{En effet, \ill{} \ill{} un homéomorphisme, il est}
\struck{\ill{}}
\note{suit une démonstration d'une douzaine de lignes, barrée de grands traits
obliques et encadrée : on y lit notamment « $f_*(\mathcal{O}_T) = \mathcal{A}$
\ill{} », « $T = \bigcup T_i$ », « $T_i = \bigcup_j T_{ij}$, où
$T_{ij} = f^{-1}(S_j) = f_i^{-1}(S_j)$ », « $T = \bigcup_{ij} f^{-1}(S_j)$ »,
et, en marge de la reprise, « \ill{} $T$ est \uncertain{quasi-compact}, on
\ill{} \ill{} $S$ affine »}

Soit $s \in S$, $f^{-1}(s)$ est réduit à un pt, donc \ill{} \ill{} \ill{}
voisinage affine $U$, lequel \ill{} \ill{} \ill{} \ill{} \ill{} voisinage de
$s$, il contient un voisinage affine $V$ de $s$. Alors \ill{} \ill{}
\[
  f^{-1}(V) = f'^{-1}(V)
\]
où $f' : U \to S$ est induit par $f$. Comme \uncertain{c'est} un morphisme de
schémas affines, et $V$ affine, $f'^{-1}(V)$ affine, O.K.

\page{18}

\emph{Proposition 2} \textit{Soit $G$ un schéma en groupes sur $S$. Pour
qu'il soit d'échelon infinitésimal $\nu$, i.e.\ $U_{\nu}(G) = G$, il faut et
il suffit que $G$ soit \emph{affine} sur $S$, et défini par une Algèbre
augmentée $\mathcal{A}$} \struck{\ill{} \ill{} \ill{}]} \textit{dont
l'Idéal $\mathcal{I}$ d'augmentation} \struck{et} \struck{tel que
$x \in \mathcal{I} \Rightarrow x^{p^{\nu}} = 0$} \add{\textit{a puissances
$p^{\nu}$-ièmes nulles}}.

D'autre part, pour $G$ quelconque, $U_{\nu}(G)$ est
\struck{\uncertain{canoniquement isomorphe}} \add{\uncertain{l'unique plus}}
\ill{} \ill{} \ill{} \ill{} \ill{} voisinage infinitésimal d'ordre $p^{\nu}$
de la section unité.
\note{« échelon » est une lecture douteuse ; la ligne barrée sous
« augmentée $\mathcal{A}$ » commence par « morphisme » ; l'insertion au-dessus
de la fin de phrase est en partie illisible}

\marginal{N.B. N'\ill{} pas \ill{} \ill{} \ill{} \uncertain{voisinage}
\ill{} \ill{} \ill{} \uncertain{section} \ill{} \ill{} $G$ \ill{} \ill{}
\ill{}}\note{dans la marge gauche, de biais}

\page{19}

\emph{Détermination de ${}_p\mathrm{Pic}(X)$, $X$ propre normal sur $k$
parfait de car.\ $p$.}

On aura
\[
  {}_p\mathrm{Pic}(X) = {}_pH^1(X, \underline{\mathcal{O}}_X^{*})
\]
Comme $X$ est \emph{réduit}, on a une suite exacte
\[
  0 \to \underline{\mathcal{O}}_X^{*} \xrightarrow{\ x \mapsto x^p\ }
  \underline{\mathcal{O}}_X^{*} \to
  \underline{\mathcal{O}}_X^{*} / \underline{\mathcal{O}}_X^{*p} \to 0
\]
d'où par la suite exacte de cohomologie, posant $A = H^0(X, \mathcal{O}_X)$
\[
  0 \to A^{*} \xrightarrow{\ x \mapsto x^p\ } A^{*} \to
  \Gamma(X, \underline{\mathcal{O}}_X^{*}/\underline{\mathcal{O}}_X^{*p})
  \to {}_pH^1(X, \underline{\mathcal{O}}_X^{*}) \to 0
\]
et comme $X$ est réduit et propre sur $k$, donc $A$ réduite finie
\struck{et séparable} sur $k$, et $k$ parfait, on a aussi
$A^{*} \simeq A^{*p}$, d'où ($X$ réduit et propre sur $k$ parfait)
\[
  {}_p\mathrm{Pic}(X) \simeq
  \Gamma(X, \underline{\mathcal{O}}_X^{*}/\underline{\mathcal{O}}_X^{*p})
\]
Considérons l'hom.\ de dérivation logarithmique
\[
  \underline{\mathcal{O}}_X^{*} \longrightarrow \underline{\Omega}^1_{X/k}
  \qquad f \longmapsto \frac{df}{f}
\]
elle est nulle sur $\underline{\mathcal{O}}_X^{*p}$, donc on a un hom.
\[
  \underline{\mathcal{O}}_X^{*}/\underline{\mathcal{O}}_X^{*p}
  \longrightarrow \underbrace{Z\underline{\Omega}^1_{X/k}}_{\text{formes fermées}}
\]
Cet hom.\ est-il injectif ? Notons que si $K$ est l'anneau des fonctions
rationnelles sur $X$, alors $f \in K$, $df = 0$ $\Longrightarrow$
$f \in K^p$ [car \ill{} \ill{} \ill{} \uncertain{ramène} au cas où $K$ est un
corps, et $K$ \struck{\ill{} \ill{}}. \ill{} \ill{} $p$-base sur
$K^p \supset k$\ldots]. Or $X$ étant \emph{normal}, si
$f \in \mathcal{O}_X$ est

\page{20}

tel que $df = 0$, \add{donc $f \in K^{*p}$} on a $f \in \mathcal{O}_X^{p}$
\add{$[\underline{\mathcal{O}}_X^{*}/\underline{\mathcal{O}}_X^{*p} \to
\underline{R}_X^{*}/\underline{R}_X^{*p}$ est injectif$]$}, donc $X$ étant
normal, on a un hom.\ \emph{injectif}
\[
  0 \to \underline{\mathcal{O}}_X^{*}/\underline{\mathcal{O}}_X^{*p}
  \longrightarrow Z\underline{\Omega}^1_{X/k}
\]
D'après Cartier, on sait que l'image est \uncertain{formée} des $\omega$ tels
que
\[
  C\omega = \omega
\]
[$C$ l'opération de Cartier], et dans le cas birationnel la réciproque est
vraie : si $\omega$ est une $1$-forme fermée telle que $C\omega = \omega$,
alors $\omega = \dfrac{df}{f}$, où $f$ est une \emph{fonction rationnelle}
qui est \add{déterminée mod} \struck{D'ailleurs, \ill{}} [sur chaque
composante de $X$] $K^{*p}$. Comme
$\underline{\mathcal{O}}_X^{*}/\underline{\mathcal{O}}_X^{*p} \to
\underline{R}_X^{*}/\underline{R}_X^{*p}$ est injectif, pour montrer que
$\omega$ provient d'une section de
$\underline{\mathcal{O}}_X^{*}/\underline{\mathcal{O}}_X^{*p}$, la question
est \emph{locale}, donc il faut prouver que si
$f \in (K = \text{anneau des fonctions de } \mathcal{O})^{*}$, telle que
$\dfrac{df}{f}$ \emph{régulière}, alors $f = g^p h$, avec
$h \in \mathcal{O}^{*}$, i.e.\ \struck{on divise} l'hom.
\[
  \underline{R}_X^{*}/\underline{\mathcal{O}}_X^{*} \longrightarrow
  \underline{\Omega}^1(\underline{K})/\underline{\Omega}^1(\underline{\mathcal{O}}_X)
\]
est \emph{injectif}. $X$ étant normal, \struck{\ill{} \ill{} \ill{}} on se
ramène au cas où $X$ est \struck{\ill{}} régulier\ldots
\note{la page s'arrête ici}

\end{document}
